\originalchapter{旋转曲面与极小曲面}
\originalsection{旋转曲面的基本形式}
旋转曲面将一个平面母线绕轴旋转, 是同时练习基本形式、曲率和测地线的贯穿算例. 令母线 $(r(s),z(s))$ 以弧长参数表示, 满足 $r>0$, $r'^2+z'^2=1$, 参数片为
\[
 X(s,\theta)=(r(s)\cos\theta,r(s)\sin\theta,z(s)).
\]
\begin{proposition}\label{prop:rotation}
取 $N=(-z'\cos\theta,-z'\sin\theta,r')$, 则
\[
 \I=\dd s^2+r^2\dd\theta^2,\qquad
 \II=(-r''z'+z''r')\dd s^2+rz'\dd\theta^2.
\]
主方向为经向与纬向, 主曲率为 $k_s=-r''z'+z''r'$, $k_\theta=z'/r$, 且 $K=-r''/r$.
\end{proposition}
\begin{proof}
$X_s=(r'\cos\theta,r'\sin\theta,z')$, $X_\theta=(-r\sin\theta,r\cos\theta,0)$, 两者正交且长度为 $1,r$. 其叉积归一化即为所给 $N$. 二阶导数取法向分量得到 $e=-r''z'+z''r'$, $f=0$, $g_2=rz'$, 所以两坐标方向为主方向.

微分单位速度条件得 $r'r''+z'z''=0$. 将 $e z'$ 展开为 $-r''z'^2+z''r'z'=-r''(z'^2+r'^2)=-r''$. 因而 $K=k_sk_\theta=-r''/r$. 此推导不除以 $z'$, 在 $z'=0$ 的点同样有效.
\end{proof}
例如球面母线 $r=R\sin(s/R),z=R\cos(s/R)$ 给出 $K=1/R^2$; 圆柱母线 $r=R,z=s$ 给出 $K=0$. 内蕴曲率只由 $r$ 的二阶变化决定, 而各主曲率仍依赖法向选择.
\begin{example}
求标准环面 $X(\theta,\varphi)=((R+a\cos\varphi)\cos\theta,(R+a\cos\varphi)\sin\theta,a\sin\varphi)$ 的面积与 Gauss 曲率, $R>a>0$.
\end{example}
\begin{solution}
直接有
\[
 \I=(R+a\cos\varphi)^2\dd\theta^2+a^2\dd\varphi^2,\qquad
 \dd A=a(R+a\cos\varphi)\dd\theta\dd\varphi.
\]
取指向管外的法向 $N=(\cos\varphi\cos\theta,\cos\varphi\sin\theta,\sin\varphi)$, 两主曲率为 $-1/a$ 和 $-\cos\varphi/(R+a\cos\varphi)$. 因而
\[
 K=\frac{\cos\varphi}{a(R+a\cos\varphi)},\qquad \Area=4\pi^2aR.
\]
环面外侧 $K>0$, 内侧 $K<0$, 上下两条纬圆 $K=0$. 积分 $K\dd A=\cos\varphi\dd\theta\dd\varphi$ 得总 Gauss 曲率0, 后面的 Gauss--Bonnet 将解释该值与形状细节无关.
\end{solution}
\begin{center}
\begin{tikzpicture}[scale=2]
\begin{axis}[width=8.6cm,height=5.2cm,view={35}{30},axis lines=none,axis equal image]
\addplot3[surf,domain=0:360,y domain=0:360,samples=25,samples y=17,shader=flat,z buffer=sort,colormap={ring}{color(0cm)=(white);color(1cm)=(structurecolor)}] ({(2+.65*cos(y))*cos(x)},{(2+.65*cos(y))*sin(x)},{.65*sin(y)});
\end{axis}
\end{tikzpicture}
\end{center}
\originalsection{直纹曲面与可展性}
\begin{definition}
由直线族组成的参数片 $X(s,t)=c(s)+td(s)$, $|d(s)|=1$, 称为直纹曲面. 正则部分若 $K=0$, 称为局部可展的直纹曲面; 这里先以曲率条件使用术语, 其局部平直性会在内蕴章节说明.
\end{definition}
\begin{proposition}\label{prop:ruled}
直纹参数片在正则点有
\[
 K=-\frac{\det(c',d,d')^2}{|(c'+td')\times d|^4}\le0.
\]
特别 $K=0$ 当且仅当 $c',d,d'$ 线性相关.
\end{proposition}
\begin{proof}
$X_s=c'+td'$, $X_t=d$, $X_{tt}=0$. 因而第二基本形式 $g_2=0$, $f=\langle d',N\rangle$. 记 $W=|X_s\times d|$, 则
\[
 f=\frac{\langle d',(c'+td')\times d\rangle}{W}
   =\frac{\det(c',d,d')}{W}.
\]
由 $EG-F^2=W^2$, 得 $K=-f^2/W^2$.
\end{proof}
平面、圆柱、圆锥的正则部分均满足零曲率. 另一例子是空间曲线的切线曲面.
\begin{example}
单位速度曲线 $c$ 有 $\kappa>0$, 考察 $X(s,t)=c(s)+tT(s)$ 在 $t>0$ 的正则部分, 求其曲率.
\end{example}
\begin{solution}
$X_s=T+t\kappa n$, $X_t=T$, 叉积为 $-t\kappa b$, 所以 $N=-b$, $E=1+t^2\kappa^2,F=1,G=1$. 有 $X_{st}=\kappa n$, 其法向分量为0, $X_{tt}=0$, 而 $X_{ss}$ 的 $b$ 分量为 $t\kappa\tau$, 故 $e=-t\kappa\tau$. 于是 $K=0$, $H=-\tau/(2t\kappa)$. $t=0$ 时叉积为0, 该参数片在回归边上奇异, 不能把公式延伸至那里.
\end{solution}
\originalsection{极小曲面的面积变分}
\begin{definition}
平均曲率 $H\equiv0$ 的曲面称为极小曲面. 名称指面积一阶变分为0, 不自动保证每一个有限区域都是全局面积最小.
\end{definition}
\begin{theorem}\label{thm:first-area}
取参数区域 $D$ 满足 $\overline D$ 紧含于光滑参数片的定义域, 并取 $f\in C_c^\infty(D)$, 考察法向变分 $X_t=X+t fN$. 当 $t$ 足够小时它保持正则, 且
\[
 \left.\frac{\dd}{\dd t}\Area(X_t(D))\right|_{t=0}
 =-2\int_D fH\dd A.
\]
在此面积按参数片积分计, 对单射片等于实际面积. 对所有这种变分一阶导数为0, 当且仅当 $H=0$.
\end{theorem}
\begin{proof}
一阶导数 $\partial_iX_t=X_i+t(f_iN+fN_i)$. 故度量变分为 $\dot g_{ij}=f(\langle N_i,X_j\rangle+\langle X_i,N_j\rangle)=-2fh_{ij}$. 对正定矩阵有 $\partial_t\det g=(\det g)\tr(g^{-1}\dot g)$, 这可由行列式的多线性或伴随矩阵公式得到. 因而
\[
 \partial_t\sqrt{\det g}\big|_0
 =\tfrac12\sqrt{\det g}\tr(-2fg^{-1}h)=-2fH\sqrt{\det g}.
\]
紧支撑使积分微分合法且边界固定. 若积分对任意 $f$ 为0, 在 $H$ 非零点附近取非负、非零紧支撑函数, 使 $H$ 在其支撑上保持同号会得到非零积分, 矛盾. 正则性在紧支撑集上由原叉积的正下界和连续性保持, 支撑外曲面不变.
\end{proof}
\begin{proposition}
图像 $z=f(u,v)$ 为极小曲面当且仅当
\begin{equation}\label{eq:minimal-graph}
 (1+f_v^2)f_{uu}-2f_uf_vf_{uv}+(1+f_u^2)f_{vv}=0.
\end{equation}
\end{proposition}
\begin{proof}
将图像的第一、第二基本形式代入命题\ref{prop:KH}, 分母为 $2(1+f_u^2+f_v^2)^{3/2}>0$, 分子恰为左侧, 所以等价. 这也等价于 $\operatorname{div}(\nabla f/\sqrt{1+|\nabla f|^2})=0$.
\end{proof}
\begin{example}
验证悬链面和螺旋面都是极小曲面, 并计算 Gauss 曲率.
\end{example}
\begin{solution}
悬链面取 $X(u,v)=(a\cosh v\cos u,a\cosh v\sin u,av)$, $a>0$. 有 $E=G=a^2\cosh^2v,F=0$. 选 $N=(\sech v\cos u,\sech v\sin u,-\tanh v)$, 二阶法向系数为 $e=-a,f=0,g_2=a$. 两主曲率为 $\mp1/(a\cosh^2v)$, 故 $H=0$, $K=-1/(a^2\cosh^4v)$.

螺旋面 $Y(u,v)=(u\cos v,u\sin v,av)$ 的 $E=1,F=0,G=u^2+a^2$. 取 $N=(a\sin v,-a\cos v,u)/\sqrt{u^2+a^2}$. 则 $e=g_2=0,f=-a/\sqrt{u^2+a^2}$, 故 $H=0$, $K=-a^2/(u^2+a^2)^2$. 两曲面都具有负 Gauss 曲率, 所以极小不等于平直.
\end{solution}
\originalsection{习题}
\begin{exercise}\label{ex:6a}
环面上取环向与管向单位主方向. 当 $R=2a$, 求环面上满足 $H=0$ 的点集, 并说明整个环面是否极小.
\end{exercise}
\begin{exercise}\label{ex:6b}
圆锥 $X(r,\theta)=(r\cos\theta,r\sin\theta,ar)$, $a>0$, $r>0$. 求两主曲率和平均曲率, 取 $X_r\times X_\theta$ 确定的法向.
\end{exercise}
\begin{exercise}\label{ex:6c}
证明没有非空紧致无边界极小嵌入曲面 $M\subset\R^3$. 提示可考察离固定原点距离最大的点, 但须写出二阶导数的符号约束.
\end{exercise}
